\documentclass[12pt, twoside]{article}
%\documentclass[12pt, twoside]{article}
\usepackage[letterpaper, margin=1in, headsep=0.2in]{geometry}
\setlength{\headheight}{0.6in}
%\usepackage[english]{babel}
\usepackage[utf8]{inputenc}
\usepackage{microtype}
\usepackage{amsmath}
\usepackage{amssymb}
%\usepackage{amsfonts}
\usepackage[nomessages]{fp} %\FPeval{\var-name}{2*sin(pi/6)}
\usepackage{siunitx} %units in math. eg 20\milli\meter
\usepackage{yhmath} % for arcs, overparenth command
\usepackage{tikz} %graphics
\usetikzlibrary{quotes, angles, arrows, arrows.meta}
\usepackage{pgfplots}
\usepgfplotslibrary{fillbetween}
\pgfplotsset{compat=1.16}
\usepackage{graphicx} %consider setting \graphicspath{{images/}}
\usepackage{parskip} %no paragraph indent
\usepackage{enumitem}
\usepackage{multicol}
\usepackage{venndiagram}

\usepackage{fancyhdr}
\pagestyle{fancy}
\fancyhf{}
\renewcommand{\headrulewidth}{0pt} % disable the underline of the header
\raggedbottom
\hfuzz=2mm %suppresses overfull box warnings

\usepackage{hyperref}
\setlength{\parskip}{0.3\baselineskip}

\fancyhead[LE]{\thepage}
\fancyhead[RO]{Markscheme}
\fancyhead[LO]{La Scuola d'Italia / Dr.\ Huson / IB Math 12 \\* 30 September 2026}

\begin{document}
\subsubsection*{1.6 Classwork: Sampling bias and survey methodologies --- Markscheme}

\emph{Calculator allowed. Marks are not printed on the student sheet; they are
recorded here. Total: 17 marks.}

\emph{\textbf{Throughout:} a technique must be named with its IB term (simple
random, convenience, systematic, quota, stratified). An explanation earns its
R1 only when it is tied to the context of the question, not a general
definition alone.}

\begin{enumerate}[itemsep=0.3cm]

%--- Problem 1: IA survey plan; population; convenience bias; stratified allocation; systematic [8 marks] ---
\item[1.] [Maximum mark: 8]

Grades 9--12: $84, 78, 68, 70$ students; total $300$. Sample size $30$.

\textbf{(a)} All the students at the high school ($300$ students).
\hfill \textbf{A1}

\emph{``The $30$ students surveyed'' or ``Grade 12 students'' earns no
mark --- that is a sample, not the population.}

\textbf{(b)} Convenience sampling \hfill \textbf{A1}

Only Grade 12 students are surveyed, and only those who choose the common
room, so the sample is not representative of all grades; older students
may have different sleep habits. \hfill \textbf{R1}

\emph{Accept any reason tied to the context: only one grade is included;
friends of Alex or students with the same routine are over-represented;
not every student has an equal chance of being chosen.}

\emph{``The sample is too small'' earns no R1 --- the sample size is the
same $30$ as the later stratified sample.}

\textbf{(c)(i)} $\dfrac{78}{300} \times 30$ \hfill \textbf{M1}

$= 7.8$ \hfill \textbf{A1}

$\approx 8$ students \hfill \textbf{AG}

\emph{Accept $\dfrac{30}{300} = \dfrac{1}{10}$ of each grade, so
$\dfrac{78}{10} = 7.8$. The value $7.8$ must be seen for the A1; writing
only ``$8$'' earns the M1 at most.}

\textbf{(c)(ii)} Grade 9: $8.4 \approx 8$; \quad Grade 11: $6.8 \approx 7$;
\quad Grade 12: $7.0 = 7$ \hfill \textbf{A1A1}

\emph{A1 for two correct, A1 for all three correct. Check: $8 + 8 + 7 + 7 =
30$.}

\emph{Numbers equal in each grade ($7.5$ or ``$7$ or $8$ each'') earn no
marks.}

\textbf{(d)} Systematic sampling \hfill \textbf{A1}

\emph{Accept ``systematic random sampling''. ``Random sampling'' alone earns
no mark.}

\newpage

%--- Problem 2: Discrete/continuous; quota; simple random; outlier as error; leading question [9 marks] ---
\item[2.] [Maximum mark: 9]

\textbf{(a)(i)} Discrete \hfill \textbf{A1}

\textbf{(a)(ii)} Continuous \hfill \textbf{A1}

\emph{Time is measured, not counted, so it is continuous even when it is
recorded to the nearest hour or half-hour.}

\textbf{(b)} Quota sampling \hfill \textbf{A1}

\emph{``Convenience'' earns no mark: the fixed number from each group is
what makes it quota sampling.}

\textbf{(c)} In quota sampling the students within each group are not
chosen at random (they are the first to arrive); in stratified sampling
they are chosen at random, and the group sizes are in proportion to the
population. \hfill \textbf{R1}

\emph{Award R1 for either the ``not random within each group'' point or the
``not in proportion to the population'' point.}

\textbf{(d)} Use a random number generator (GDC or spreadsheet) to produce
random whole numbers from $1$ to $300$ \hfill \textbf{M1}

select the students with those numbers, ignoring repeats, until $30$
different students are chosen \hfill \textbf{A1}

\emph{Accept drawing $30$ numbers from a hat containing $1$ to $300$. M1 for
a random method applied to the numbered list; A1 for $30$ distinct
students (no repeats). ``Pick $30$ names from the list'' with no random
method earns no marks.}

\emph{\textbf{GDC:} \texttt{randInt(1,300)} repeated, or
\texttt{randSamp} on the list $1$ to $300$ without replacement.}

\textbf{(e)} Check the value with the student, or remove it from the data
set \hfill \textbf{A1}

$19$ hours of sleep in one night is unrealistic, so the value is most likely
a recording error (for example $9$ hours) \hfill \textbf{R1}

\emph{Accept ``it is an outlier and is probably an error''. The R1 needs the
reason that the value is implausible or an error; ``it is an outlier'' alone
earns the A1 only. ``Keep it'' earns A0, even with the reason that some
outliers are valid data.}

\textbf{(f)} It is a leading question: its wording suggests the answer
``yes'', so students may agree rather than give their honest opinion.
\hfill \textbf{R1}

\emph{Accept ``biased wording'', ``it pushes students towards one answer''.
``Students may lie'' alone earns no mark --- the R1 is for identifying that
the question itself causes the bias.}

\end{enumerate}
\end{document}
